
\documentclass{article}
\usepackage{colortbl}
\usepackage{makecell}
\usepackage{multirow}
\usepackage{supertabular}

\begin{document}

\newcounter{utterance}

\centering \large Interaction Transcript for game `cladder', experiment `full\_v1.5\_default', episode 4837 with Qwen3.6{-}27B.
\vspace{24pt}

{ \footnotesize  \setcounter{utterance}{1}
\setlength{\tabcolsep}{0pt}
\begin{supertabular}{c@{$\;$}|p{.15\linewidth}@{}p{.15\linewidth}p{.15\linewidth}p{.15\linewidth}p{.15\linewidth}p{.15\linewidth}}
    \# & \multicolumn{2}{c}{Player} && \multicolumn{2}{c}{Game Master} \\
    \hline

    \theutterance \stepcounter{utterance}  
    & & & \multicolumn{4}{p{0.6\linewidth}}{
        \cellcolor[rgb]{0.9,0.9,0.9}{
            \makecell[{{p{\linewidth}}}]{
                \texttt{\tiny{[P1$\langle$GM]}}
                \texttt{You are an expert in causal inference. The following question is not a typical commonsense query, but rather a meticulously designed question created by a professor specializing in causal inference, intended to assess the students' mastery of the course content.} \\
\texttt{Imagine a self{-}contained, hypothetical world with only the following conditions, and without any unmentioned factors or causal relationships: Ability to speak english has a direct effect on the forest. The smoker has a direct effect on the forest. The overall probability of speaking english is 93\%. For people who do not speak english, the probability of the forest on fire is 79\%. For people who speak english, the probability of the forest on fire is 65\%. Is the forest on fire more likely than the forest not on fire overall?} \\
\texttt{Start your answer with 'Yes' or 'No', followed by additional reasoning or evidence to support your explanation.} \\
            }
        }
    }
     \\ \\

    \theutterance \stepcounter{utterance}  
    & \multicolumn{4}{p{0.6\linewidth}}{
        \cellcolor[rgb]{0.561,0.737,0.561}{
            \makecell[{{p{\linewidth}}}]{
                \texttt{\tiny{[P1$\rangle$GM]}}
                \texttt{No} \\
\\ 
\texttt{To determine whether the forest is on fire more likely than not overall, we must calculate the total probability of the forest being on fire, denoted as \$P(F)\$.} \\
\\ 
\texttt{According to the Law of Total Probability, we can express \$P(F)\$ by conditioning on the variable "Ability to speak English" (\$E\$):} \\
\\ 
\texttt{\$\$P(F) = P(F | E) \textbackslash{}cdot P(E) + P(F | \textbackslash{}neg E) \textbackslash{}cdot P(\textbackslash{}neg E)\$\$} \\
\\ 
\texttt{From the problem description, we have the following values:} \\
\texttt{1.  The overall probability of speaking English: \$P(E) = 0.93\$} \\
\texttt{2.  Therefore, the probability of not speaking English: \$P(\textbackslash{}neg E) = 1 {-} 0.93 = 0.07\$} \\
\texttt{3.  The probability of the forest on fire given the person does not speak English: \$P(F | \textbackslash{}neg E) = 0.79\$} \\
\texttt{4.  The probability of the forest on fire given the person speaks English: \$P(F | E) = 0.65\$} \\
\\ 
\texttt{Note: The information regarding "The smoker has a direct effect on the forest" is extraneous to this specific calculation because no probabilities or conditional relationships involving smokers are provided. In a self{-}contained hypothetical world defined by the given conditions, we rely only on the quantified variables provided.} \\
\\ 
\texttt{Now, we substitute the values into the equation} \\
            }
        }
    }
    & & \\ \\

    \theutterance \stepcounter{utterance}  
    & & & \multicolumn{2}{p{0.3\linewidth}}{
        \cellcolor[rgb]{0.9,0.9,0.9}{
            \makecell[{{p{\linewidth}}}]{
                \texttt{\tiny{[GM$|$GM]}}
                \texttt{No} \\
            }
        }
    }
    & & \\ \\

    \theutterance \stepcounter{utterance}  
    & & & \multicolumn{2}{p{0.3\linewidth}}{
        \cellcolor[rgb]{0.9,0.9,0.9}{
            \makecell[{{p{\linewidth}}}]{
                \texttt{\tiny{[GM$|$GM]}}
                \texttt{yes} \\
            }
        }
    }
    & & \\ \\

    \theutterance \stepcounter{utterance}  
    & & & \multicolumn{2}{p{0.3\linewidth}}{
        \cellcolor[rgb]{0.9,0.9,0.9}{
            \makecell[{{p{\linewidth}}}]{
                \texttt{\tiny{[GM$|$GM]}}
                \texttt{game\_result = LOSE} \\
            }
        }
    }
    & & \\ \\

\end{supertabular}
}

\end{document}
